% GENERATED FILE. Edit canonical lessons, then run npm run generate:books.
% canonical-source-hash: fnv1a64:c909b594341afb2e
% canonical-lessons: PT-C330-descascar, PT-C330-picar, PT-C330-ralar, PT-C330-desfiar, PT-C330-tigela

\chapter{Preparar os legumes}
\label{ch:pt-preparar-os-legumes}
\hlchaptermodality{330}

\begin{chapteropening}
\textbf{By the end of this chapter:} \emph{I can describe getting food ready: peeling, chopping, grating and shredding, and the bowl it goes in.}

Complete the last lesson of chapter 330: I can describe getting food ready: peeling, chopping, grating and shredding, and the bowl it goes in.
\end{chapteropening}

\section[descascar]{descascar — to peel}

\label{lesson:PT-C330-descascar}



\begin{quote}
Before the new one: say the Portuguese for to conquer, then the Portuguese for the revolution.
\end{quote}

\subsection*{You'll want to know: descascar}
\textbf{descascar} — \textquotedblleft{}to peel\textquotedblright{}.

\begin{quote}
A minha mãe descasca as batatas.
\end{quote}

\begin{quote}
\emph{My mother peels the potatoes.}
\end{quote}

Built on \emph{a casca}, ‘the peel, the shell’.

\begin{cousinweb}
\emph{descascar} is built from \emph{des-} and \emph{casca} \textquotedblleft{}peel, shell, bark\textquotedblright{}; EtymDB derives \emph{casca} from the verb \emph{cascar} and goes no further, so the deeper origin is not settled.
\end{cousinweb}

\begin{grammarlens}[title={What you've built}]
One more word for telling, explaining and describing at B1.
\end{grammarlens}

\subsection*{Your turn}
\begin{itemize}
\raggedright
  \item \emph{Say it:} \emph{descascar}
  \item \emph{Say it:} \emph{descascar}, once more
  \item \emph{Say it:} a sentence of your own with \emph{descascar}
  \item \emph{Recall:} say the Portuguese for to conquer, then the Portuguese for the revolution, then say \emph{descascar} again
\end{itemize}

\begin{tcolorbox}[breakable,colback=teal!4,colframe=teal!35!black,title={Before you move on}]
What does descascar mean? (\textquotedblleft{}To peel\textquotedblright{}.) Say it once more.
\end{tcolorbox}
\section[picar]{picar — to chop finely, to mince; to sting}

\label{lesson:PT-C330-picar}



\begin{quote}
Before the new one: say the Portuguese for the revolution, then the Portuguese for to peel.
\end{quote}

\subsection*{You'll want to know: picar}
\textbf{picar} — \textquotedblleft{}to chop finely, to mince; to sting\textquotedblright{}.

\begin{quote}
O cozinheiro picou a cebola e o alho.
\end{quote}

\begin{quote}
\emph{The cook chopped the onion and garlic finely.}
\end{quote}

Also of insects: \emph{a abelha picou-me}. \emph{Carne picada} — mince.

\begin{cousinweb}
\emph{picar} is inherited from Vulgar Latin \emph{piccāre} \textquotedblleft{}to prick, to puncture\textquotedblright{} (EtymDB).
\end{cousinweb}

\begin{grammarlens}[title={What you've built}]
One more word for telling, explaining and describing at B1.
\end{grammarlens}

\subsection*{Your turn}
\begin{itemize}
\raggedright
  \item \emph{Say it:} \emph{picar}
  \item \emph{Say it:} \emph{picar}, once more
  \item \emph{Say it:} a sentence of your own with \emph{picar}
  \item \emph{Recall:} say the Portuguese for the revolution, then the Portuguese for to peel, then say \emph{picar} again
\end{itemize}

\begin{tcolorbox}[breakable,colback=teal!4,colframe=teal!35!black,title={Before you move on}]
What does picar mean? (\textquotedblleft{}To chop finely, to mince; to sting\textquotedblright{}.) Say it once more.
\end{tcolorbox}
\section[ralar]{ralar — to grate}

\label{lesson:PT-C330-ralar}



\begin{quote}
Before the new one: say the Portuguese for to peel, then the Portuguese for to chop finely.
\end{quote}

\subsection*{You'll want to know: ralar}
\textbf{ralar} — \textquotedblleft{}to grate\textquotedblright{}.

\begin{quote}
Ela ralou o queijo para a massa.
\end{quote}

\begin{quote}
\emph{She grated the cheese for the pasta.}
\end{quote}

\emph{Queijo ralado} — grated cheese. Informally, \emph{ralar-se} is ‘to worry’.

\begin{cousinweb}
\emph{ralar} is usually linked to Latin \emph{rādere} \textquotedblleft{}to scrape\textquotedblright{} (Lewis \& Short), but the path is not settled; EtymDB has the word with no source.
\end{cousinweb}

\begin{grammarlens}[title={What you've built}]
One more word for telling, explaining and describing at B1.
\end{grammarlens}

\subsection*{Your turn}
\begin{itemize}
\raggedright
  \item \emph{Say it:} \emph{ralar}
  \item \emph{Say it:} \emph{ralar}, once more
  \item \emph{Say it:} a sentence of your own with \emph{ralar}
  \item \emph{Recall:} say the Portuguese for to peel, then the Portuguese for to chop finely, then say \emph{ralar} again
\end{itemize}

\begin{tcolorbox}[breakable,colback=teal!4,colframe=teal!35!black,title={Before you move on}]
What does ralar mean? (\textquotedblleft{}To grate\textquotedblright{}.) Say it once more.
\end{tcolorbox}
\section[desfiar]{desfiar — to shred (into strands); to fray}

\label{lesson:PT-C330-desfiar}



\begin{quote}
Before the new one: say the Portuguese for to chop finely, then the Portuguese for to grate.
\end{quote}

\subsection*{You'll want to know: desfiar}
\textbf{desfiar} — \textquotedblleft{}to shred (into strands); to fray\textquotedblright{}.

\begin{quote}
O meu pai desfia o bacalhau à mão.
\end{quote}

\begin{quote}
\emph{My father shreds the salt cod by hand.}
\end{quote}

Also a fabric that frays: \emph{a camisola desfiou}.

\begin{cousinweb}
\emph{desfiar} is built from \emph{des-} and \emph{fio} \textquotedblleft{}thread\textquotedblright{}, which is inherited from Latin \emph{fīlum} \textquotedblleft{}a thread\textquotedblright{} (Lewis \& Short): to pull apart thread by thread.
\end{cousinweb}

\begin{grammarlens}[title={What you've built}]
One more word for telling, explaining and describing at B1.
\end{grammarlens}

\subsection*{Your turn}
\begin{itemize}
\raggedright
  \item \emph{Say it:} \emph{desfiar}
  \item \emph{Say it:} \emph{desfiar}, once more
  \item \emph{Say it:} a sentence of your own with \emph{desfiar}
  \item \emph{Recall:} say the Portuguese for to chop finely, then the Portuguese for to grate, then say \emph{desfiar} again
\end{itemize}

\begin{tcolorbox}[breakable,colback=teal!4,colframe=teal!35!black,title={Before you move on}]
What does desfiar mean? (\textquotedblleft{}To shred (into strands); to fray\textquotedblright{}.) Say it once more.
\end{tcolorbox}
\section[a tigela]{a tigela — the bowl}

\label{lesson:PT-C330-tigela}



\begin{quote}
Before the new one: say the Portuguese for to grate, then the Portuguese for to shred.
\end{quote}

\subsection*{You'll want to know: a tigela}
\textbf{a tigela} — \textquotedblleft{}the bowl\textquotedblright{}.

\begin{quote}
Ele pôs a fruta numa tigela.
\end{quote}

\begin{quote}
\emph{He put the fruit in a bowl.}
\end{quote}

\emph{Uma tigela de sopa} — a bowl of soup.

\begin{cousinweb}
\emph{tigela} is inherited from Vulgar Latin \emph{tēgella} (EtymDB), a little \emph{tēgula} \textquotedblleft{}a roof tile\textquotedblright{}: a curved piece of baked clay.
\end{cousinweb}

\begin{grammarlens}[title={What you've built}]
One more word for telling, explaining and describing at B1.
\end{grammarlens}

\subsection*{Your turn}
\begin{itemize}
\raggedright
  \item \emph{Say it:} \emph{a tigela}
  \item \emph{Say it:} \emph{a tigela}, once more
  \item \emph{Say it:} a sentence of your own with \emph{a tigela}
  \item \emph{Recall:} say the Portuguese for to grate, then the Portuguese for to shred, then say \emph{a tigela} again
\end{itemize}

\begin{tcolorbox}[breakable,colback=teal!4,colframe=teal!35!black,title={Before you move on}]
What does a tigela mean? (\textquotedblleft{}The bowl\textquotedblright{}.) Say it once more.
\end{tcolorbox}
